\section{Equipotência}

No início do Capítulo~3, introduzimos a noção de \emph{equipotência} $\approx$ (definição \ref{def:equipotencia}) como forma de iniciar a discussão sobre números naturais.
Vimos ainda que tal relação constitui uma relação de equivalência no universo dos conjuntos (Proposição \ref{prop:equipotencia-equivalencia}).

Nesta seção, estudaremos com mais profundidade aquela noção.

\begin{definition}\label{def:menorIgual}
    Dizemos que a cardinalidade de um conjunto $A$ é menor ou igual à cardinalidade de um conjunto $B$ (escrevemos $A\preceq B$ --- teremos uma notação mais usual no futuro) se existe uma função injetora de $A$ em $B$.

    Dizemos que a cardinalidade de um conjunto $A$ é menor do que a cardinalidade de um conjunto $B$ (escrevemos $A\prec B$) se $A\preceq B$ e $A\not\approx B$.
\end{definition}

O lema abaixo, cuja demonstração é deixada como exercício ao leitor, diz que $\preceq$ é uma pré-ordem no universo dos conjuntos.

\begin{lemma}\label{lem:pre-ordem}
    Para quaisquer conjuntos $A$, $B$ e $C$, temos:
    \begin{enumerate}[label=(\roman*)]
        \item $A\preceq A$;
        \item se $A\preceq B$ e $B\preceq C$, então $A\preceq C$;
    \end{enumerate}
\end{lemma}
\subsection*{Exercícios}
\begin{exercise}
    Prove o Lema \ref{lem:pre-ordem}.
\end{exercise}
